\clearpage
\scanpage{5}

C) \uline{\unsure{Mèthodes|Méthodes} \inserted{de construction} projectives}

\S\,14\quad Familles \unsure{limitées}.

\S\,15\quad \inserted{Foncteurs et} Schémas de Hilbert. (+ appl. aux quotients).

\S\,16\quad \inserted{Foncteurs et} Schémas de Picard.
\margin{(\unsure{Ths|Thes} d'existence projectifs \unsure{exclusivement}\deleted{\unsure{)}}; également le cas de \unsure{gpe|gpes} \unsure{structural|structuraux} affine \unsure{commutatif} plus général que $\unsure{\mathrm G_{\mathrm m}|\mathrm G_m}$\ldots)}

\S\,17\quad Variantes de Picard.

\S\,1\deleted{\unsure{7}}8\quad Applications aux courbes

\S\,1\deleted{\unsure{8}}\unsure{9|8}\quad Formalisme des normes

\S\,\deleted{\unsure{19}}\unsure{20}\quad Invariants multiprojectifs de Mumford.

D) \uline{Méthodes \inserted{de construction} mixtes}

\S\,2\deleted{0}1\quad Cas de représentabilité de $\underline{\mathrm{Pic}}$, $\underline{\mathrm{Pic}}^{0}$, $\underline{\deleted{\unsure{\mathrm P}}\mathrm{Pic}}_{X}^{\unsure{\tau|\mathrm t}}$ \ldots

(y compris le cas de Mumford utilisant \S\,\deleted{\unsure{19}}20).

\S\,21\quad Cas de représentabilité des foncteurs $\prod_{\deleted{\unsure{X}}Y/S}X/Y$ etc.

(\unsure{ths} du type Matsumura).

\noindent\rule{\linewidth}{.3pt}

NB.\quad On suppose connues \unsure{toutes} les \unsure{gén.} \unsure{fonctorielles} (topologies, descente faisceaux, quotients) ; à rappeler au besoin dans $\unsure{0|O}_{\underline{\overline{\mathrm V}}}$.

\noindent\rule{\linewidth}{.3pt}

\uline{Manquent} :

\reversemarginpar\margin{\deleted{Il \unsure{manquera}}}\margin{dans Chap. sur schémas en groupes}
$\left\{\vphantom{\begin{matrix}a\\a\end{matrix}}\right.$ les \unsure{ths} relatifs aux VA

\margin{dans Chapitre sur Picard}
\noindent$\left\{\begin{minipage}{.83\linewidth}les \unsure{ths g\textsuperscript{enx}|ths génx} \unsure{pointus|partiels} sur Picard, et l'étude qualitative générale des schémas de Picard, (y compris \unsure{critère} de \unsure{lissité} de Mumford\ldots)

(?) \deleted{\unsure{Relation} entre}\inserted{Étude du} Picard d'un schéma projectif en termes d'un pinceau de sections hyperplanes.

Les th. sur Picard genre ``\unsure{certaines|certines|sections} \unsure{équivalences|équivalentes|équivalence}''.\end{minipage}\right.$

\margin{dans Chap. sur $\unsure{\pi_1|\pi}$}
\noindent$\left\{\begin{minipage}{.83\linewidth}Les formulations de descente à la van Kampen etc, en termes de groupe fondamental\end{minipage}\right.$
\normalmarginpar
% archive sheet number: 4
